\documentclass[12pt]{article}

%%%%%%%%%%%%%%%%%%%%%%%%%%%%%%%%%%%%%%%%%%%%%%%%%%%%%%%%%%
\usepackage{graphicx} 
\usepackage{fancyhdr}
\usepackage[misc]{ifsym}
\usepackage{enumitem}
\usepackage{tikz,tcolorbox}
\usepackage{amsmath}
\usepackage{multirow}
\usepackage{natbib}
\usepackage{algorithm}
\usepackage{algpseudocode}
\usetikzlibrary{calc}
\usepackage{amsfonts}
\usepackage[flushleft]{threeparttable} % Required for tablenotes
\usepackage{pgfplots}
\pgfplotsset{width=10cm}
\usepackage[margin=0.8in]{geometry}
\usepackage{amsthm}
\usepackage{mathtools}
\usepackage{wasysym}
\usepackage{listings}
\usepackage{amssymb}
\usepackage{exercise}
\usepackage{appendix}
\usepackage{booktabs} % For better table formatting
\usepackage{geometry} % To adjust page margins
\geometry{margin=1in}
\usepackage{titlesec}
\usepackage{titlesec}
% The [runin] shape forces the text to start on the same line
\titleformat{\subsubsection}[runin]
  {\normalfont\normalsize\bfseries}{\thesubsubsection}{1em}{}[:\quad]
\newcommand{\appendixsection}[1]{
  \section*{#1} % A section without numbering
  \addcontentsline{toc}{section}{#1} % Add it to the table of contents
}
\titleformat{\appendixsection}
  {\normalfont\HUGE\bfseries}{\thesection.}{1em}{}



\newcommand{\piercedheart}{%
  \relbar
  \mathrel{\mspace{-3.5mu}}
  \mathrel{\loweredheartsuit}
  \mathrel{\mspace{-7mu}}
  \rightarrow
}

\makeatletter
\newcommand{\loweredheartsuit}{\mathpalette\loweredheartsuit@\relax}
\newcommand{\loweredheartsuit@}[2]{%
  \vcenter{%
    \sbox\z@{$\m@th#1\mspace{3mu}$}%
    \vskip\wd\z@
    \hbox{$\m@th#1\heartsuit$}%
  }%
}
\makeatother


\lstset{
	basicstyle=\selectfont\ttfamily
}

\usepackage{hyperref}
 \hypersetup{
     colorlinks=true,    
     urlcolor=blue,
     citecolor=red, 
     linkcolor=blue,
 }

\usepackage{epigraph}
\setlength \epigraphwidth {12cm}
\setlength \epigraphrule {0pt}
\AtBeginDocument{\renewcommand {\epigraphflush}{center}}



\setlength{\headheight}{15pt}

\pgfplotsset{compat=1.18}


\newtheorem*{claim}{Claim}
\newtheorem*{definition}{Definition}

%%%%%%%%%%%%%%%%%%%%%%%%%%%%%%%%%%%%%%%%%%%%%%%%%%%%%%%%%%
\begin{document}

\pagestyle{fancy}
\fancyhead{} % clear all header fields
\fancyhead[L]{\textbf{IE 514 Network Flows}}


%%%%%%%%%%%%%%%%%%%%%%%%%%%%%%%%%%%%%%%%%%%%%%%%%%%%%%%%%%
\input{titlepage}
%%%%%%%%%%%%%%%%%%%%%%%%%%%%%%%%%%%%%%%%%%%%%%%%%%%%%%%%%%

\section{Introduction}
The Degree-Constrained Minimum Spanning Tree (DCMST) problem seeks to find a spanning tree of minimum total weight such that the degree of each vertex does not exceed a predefined limit. As a classic NP-hard combinatorial optimization problem, it finds extensive application in telecommunications, electrical grid layout, and causal network design. Due to its computational complexity, the literature is roughly divided into exact mathematical programming techniques, relaxation-based bounds, and metaheuristics.

\section{Formulations}

Let $V$ be the set of vertices, $E$ the set of undirected edges, and $\delta(v)$ the set of edges incident to vertex $v \in V$. 
The set of all valid, unconstrained spanning trees of $G=(V,E)$ is denoted by $\mathcal{T}$. 
The cost of an edge $e \in E$ is $w_e \in \mathbb{R}$, and the maximum degree allowed for vertex $v \in V$ is $d_v \in \mathbb{Z}^+$. 
Binary decision variables $x_e \in \{0,1\}$ indicate whether edge $e$ is selected ($x_e=1$) or not.

\subsection{Non‑compact Formulation} \label{sec:non-compact}
The non‑compact formulation models the spanning tree directly with undirected edge variables and explicitly forbids all cycles through an exponential family of subset constraints.
\begin{align}
\min \quad & \sum_{e \in E} w_e \, x_e \label{nc:obj} \\
\text{s.t.} \quad & \sum_{e \in \delta(v)} x_e \le d_v && \forall v\in V \label{nc:degree} \\
& \sum_{e \in E} x_e = |V| - 1 \label{nc:edgecount} \\
& \sum_{e \in E(S)} x_e \le |S| - 1 && \forall S \subset V,\; |S|\ge 2 \label{noncompact} \\
& x_e \in \{0,1\} && \forall e\in E \label{nc:binary}
\end{align}
Here $E(S)$ denotes the set of edges with both endpoints in $S$.  
Constraint \eqref{nc:degree} enforces the maximum degree, \eqref{nc:edgecount} fixes the correct number of edges, and the subtour elimination constraints \eqref{noncompact} guarantee that no proper subset of vertices induces a cycle.

\subsection{Compact Flow Formulations}
Compact formulations avoid the exponentially many subtour elimination constraints \eqref{noncompact} by introducing auxiliary variables (flows or node potentials) and working with directed arcs.  
Choose an arbitrary root node $r \in V$ and replace each undirected edge $e = \{i,j\} \in E$ with two directed arcs $(i,j)$ and $(j,i)$.  
Let $A = \{ (i,j), (j,i) : \{i,j\} \in E \}$ be the set of all directed arcs.  
Directed binary variables $y_{ij} \in \{0,1\}$ indicate whether arc $(i,j) \in A$ is included in the arborescence rooted at $r$.  
The undirected edge variables are recovered by the logical relation
\[
x_{\{i,j\}} = y_{ij} + y_{ji} \qquad \forall \{i,j\}\in E.
\]

\subsubsection{Multi‑commodity flow} \label{sec:multi}
For every non‑root node $k \in V \setminus \{r\}$, one unit of flow of commodity $k$ is sent from $r$ to $k$.  
Continuous flow variables $f_{ij}^k \ge 0$ denote the amount of commodity $k$ on arc $(i,j)\in A$.  
The formulation reads:
\begin{align}
\sum_{j: (i,j) \in A} f_{ij}^k \;-\; \sum_{j: (j,i) \in A} f_{ji}^k &=
\begin{cases}
 1, & i = r \\
-1, & i = k \\
 0, & \text{otherwise}
\end{cases}
&& \forall k \in V\setminus\{r\},\; \forall i\in V \label{mf:flowcons} \\
f_{ij}^k &\le y_{ij}
&& \forall (i,j)\in A,\; \forall k\in V\setminus\{r\} \label{mf:capacity}
\end{align}
Constraints \eqref{mf:flowcons} enforce flow conservation for each commodity, and \eqref{mf:capacity} ensure that flow can only use selected arcs.

\subsubsection{Single‑commodity flow}\label{sec:integer}
Aggregating all commodities gives a single integer flow.  
Send $|V|-1$ units from the root $r$; every other node consumes exactly one unit.  
Continuous flow variables $f_{ij}^k \ge 0$ for each $(i,j)\in A$:
\begin{align}
\sum_{j: (i,j) \in A} f_{ij} \;-\; \sum_{j: (j,i) \in A} f_{ji} &=
\begin{cases}
|V|-1, & i = r \\
-1,    & i \in V\setminus\{r\}
\end{cases}
&& \forall i\in V \label{sf:flowcons} \\
f_{ij} &\le (|V|-1)\, y_{ij}
&& \forall (i,j)\in A \label{sf:capacity}
\end{align}
Constraint \eqref{sf:flowcons} is the aggregated flow conservation; \eqref{sf:capacity} links flow to the arborescence arcs.

\subsubsection{MTZ potential constraints}
This formulation uses node potential variables $u_i \in \mathbb{Z}^+$ instead of explicit flows.  
Every non‑root node must have exactly one incoming arc:
\begin{align}
\sum_{i: (i,j) \in A} y_{ij} &= 1 \qquad \forall j \in V\setminus\{r\} \label{mtz:indegree} \\
u_i - u_j + |V|\, y_{ij} &\le |V|-1 \qquad \forall (i,j)\in A,\; j \neq r \label{mtz:potential}
\end{align}
Inequality \eqref{mtz:potential} together with \eqref{mtz:indegree} prevents cycles and ensures connectivity from the root.

\section{Literature Review}
The computational difficulty of the Degree-Constrained Minimum Spanning Tree (DCMST) is strictly governed by the degree parameter, typically denoted as $b_i$. As the allowed maximum degree increases, the problem relaxes toward the unconstrained Minimum Spanning Tree (MST), which is solvable in polynomial time. Conversely, tight bounds (e.g., $b_i = 2$) severely restrict the feasible space, mimicking the NP-hard Hamiltonian path problem.

\subsection{Exact Methods}

\subsubsection{Standard and Primal Branch-and-Cut Algorithms}
\citet{Caccetta2001} established early benchmarking standards using a branch-and-cut methodology. Their computational testing on random networks ranging from 100 to 800 vertices explicitly demonstrated the inverse relationship between the uniform degree bound $b_i$ and computational time. For instance, on 500-node networks, relaxing the maximum degree constraint from $b_i = 3$ to $b_i = 6$ resulted in an exponential drop in required CPU time. Consequently, their largest networks (800 nodes) were only practically solvable under looser degree constraints. 

\vspace{0.5cm}
\noindent \citet{Behle2007} expanded the exact methodology by introducing a primal branch-and-cut algorithm. Standard branch-and-cut approaches operate on the dual space, but the primal method continuously maintains a feasible DCMST solution throughout execution. Benchmarked on instances scaling up to 300 nodes, their primal separation routines yielded exceptionally strong bounds for tightly constrained networks, proving computationally competitive against the established heuristics of the time.

\subsubsection{Branch-and-Cut-and-Price Algorithms}
The current state-of-the-art exact benchmarking is defined by \citet{Bicalho2016}. Recognizing that uniformly bounded networks ($b_i \ge 3$) are computationally trivial for modern solvers, they explicitly omitted them from their benchmarks. Instead, they evaluated their branch-and-cut-and-price (BCP) algorithms entirely on heavily constrained datasets, revealing distinct performance tiers based on constraint slack:

\vspace{0.5cm}
\noindent\textit{Tight Bounds (DE/DR datasets):} For networks where vertex bounds were tightly restricted to $b_i \in \{1, 2, 3\}$, their combined Column Generation and Branch-and-Cut (CG-BC) algorithm dominated standard mathematical programming, successfully solving dense instances of up to 500 nodes.

\vspace{0.5cm}
\noindent\textit{Zero-Slack Bounds (LH datasets):} \citet{Bicalho2016} introduced a notoriously difficult benchmarking dataset where the sum of all vertex degree limits exactly equals $2(n-1)$. Because every valid spanning tree consumes the entire degree allowance, these instances possess absolute zero slack, rendering standard local-search edge-swap heuristics practically useless. Despite this extreme computational difficulty, their exact framework scaled to solve zero-slack Euclidean instances up to 800 nodes, and random-cost instances up to 600 nodes, within a 4-hour limit. \textit{They claim harder the better? suspicious}

\subsection{Relaxations and Heuristics}
 Lagrangian relaxation serves as a primary tool for establishing strong dual bounds. \citet{Volgenant1989} pioneered the approach of relaxing the degree constraints into the objective function. Utilizing the subgradient optimization principles outlined in \citet{Fisher1981}, this method iteratively improves the lower bound, making it an excellent warm-start mechanism for exact branch-and-bound trees. While direct polyhedral studies of the DCMST are rare, adjacent problems provide necessary valid inequalities. For instance, \citet{Gamarnik2005} investigated the Traveling Salesperson Problem (TSP) in cubic graphs; because the TSP is equivalent to a strict maximum-degree-2 spanning tree, their polyhedral analysis regarding strong upper bounds and subtour elimination is structurally applicable to the DCMST.

\vspace{0.5cm}
\noindent For large-scale instances exceeding 1,000 nodes where exact methods experience combinatorial explosion, metaheuristics remain highly competitive. \citet{Knowles1999} utilized a steady-state Genetic Algorithm (GA) with specialized edge-set encoding, while later Ant Colony Optimization (ACO) variants provided robust exploration capabilities. Although these methods cannot guarantee optimality or provide the strong lower bounds characteristic of exact pricing algorithms, they remain standard baselines for evaluating computational time limits.

\subsection{Position of the Present Work}
Among exact approaches, the branch-and-cut-and-price algorithm of \citet{Bicalho2016} currently defines the state of the art, solving up to 500 nodes to optimality through the combined use of column generation and dynamically generated cutting planes. 
In contrast, the branch-and-price algorithm developed in this study relies exclusively on a Dantzig-Wolfe decomposition in which the master problem is a convex combination of spanning trees and the pricing subproblem remains an unconstrained minimum spanning tree. No cutting planes are added at any stage; branching is performed directly on the original edge variables. Furthermore, because every generated column represents a complete spanning tree, we can efficiently filter out trees that violate these branching decisions, providing warm-started master problem at every child node. This design actively isolates the pure bounding power of column generation and ensures the pricing subproblem can be solved trivially using Kruskal's algorithm, bypassing the complex separation routines inherent to modern BCP frameworks.

\section{Branch and Price}

\subsection{Master Problem (MP)}

For each spanning tree $t \in \mathcal{T}$, let $a_{ev}^t \in \{0,1\}$ indicate whether edge $e \in \delta(v)$ is present in $t$, so that the degree of vertex $v$ in $t$ is $D_v^t = \sum_{e \in \delta(v)} a_{ev}^t$, and the total weight of $t$ is $W_t = \sum_{e \in E} w_e x_e^t$. 
In the master problem, binary variables $\lambda_t \in \{0,1\}$ indicate whether spanning tree $t$ is selected (relaxed to $\lambda_t \geq 0$ in the restricted master problem for column generation).
\begin{align}
\min \quad & \sum_{t \in \mathcal{T}} W_t \lambda_t \\
\text{s.t.} \quad & \sum_{t \in \mathcal{T}} D_v^t \lambda_t \leq d_v & \forall v \in V \quad & ( \pi_v \leq 0) \\
& \sum_{t \in \mathcal{T}} \lambda_t = 1 & & ( \mu) \\
& \lambda_t \geq 0 & \forall t \in \mathcal{T}
\end{align}

\noindent The Dual of the Master Problem(DMP) is formulated as:
\begin{align}
\max \quad & \sum_{v \in V} d_v \pi_v + \mu \\
\text{s.t.} \quad & \sum_{v \in V} D_v^t \pi_v + \mu \leq W_t & \forall t \in \mathcal{T} \label{eq:dual_constraint} \\
& \pi_v \leq 0 & \forall v \in V \\
& \mu \in \mathbb{R}
\end{align}

\noindent 
The constraints in the Dual Master Problem \eqref{eq:dual_constraint} require that for every valid spanning tree $t \in \mathcal{T}$, the condition $W_t - \sum_{v \in V} D_v^t \pi_v - \mu \geq 0$ must hold. This expression is exactly the {reduced cost} ($\bar{c}$) of a tree in the primal space. The column generation Subproblem searches for a tree $t$ that violates this dual constraint (i.e., finding a tree with a negative reduced cost) to add to the Master Problem.

\subsection{Pricing Subproblem (SP)}

The Subproblem generates a new spanning tree $t$ to enter the Master Problem by minimizing the reduced cost. Let $\pi_v$ be the dual variable associated with the degree constraint of vertex $v$ in the Master Problem, and $\mu$ be the dual variable for the convexity constraint. The reduced cost $\bar{c}$ of a new tree represented by edge variables $x_e$ is:
\begin{equation*}
\bar{c} = \sum_{e \in E} w_e x_e - \sum_{v \in V} \pi_v \left( \sum_{e \in \delta(v)} x_e \right) - \mu
\end{equation*}

\noindent Every edge $e = (u,v)$ contributes exactly once to the degree of $u$ and once to the degree of $v$. Therefore, we can rewrite the objective function of the subproblem by defining {modified edge weights}
$\bar{w}_{uv} = w_{uv} - \pi_u - \pi_v$. Then {pricing subproblem formulation as follows.}
\begin{align}
\min \quad & \sum_{e=(u,v) \in E} \bar{w}_{uv} x_{uv} \\
\text{s.t.} \quad & \sum_{e \in E} x_e = |V| - 1 \\
& \sum_{e \in E(S)} x_e \leq |S| - 1 & \forall S \subset V, |S| \geq 2 \\
& x_e \in \{0,1\} & \forall e \in E
\end{align}

\noindent The subproblem is simply an {Unconstrained Minimum Spanning Tree} problem. Because the tree constraints are naturally integral, the subproblem does not require an IP solver. It can be solved optimally and efficiently in $O(|E| \log |V|)$ time using {Kruskal's Algorithm} over the graph with the updated edge weights $\bar{w}_{uv}$. If the weight of the resulting MST minus $\mu$ is strictly less than 0 (i.e., negative reduced cost), the new tree is added as a column to the Restricted Master Problem.

\subsection{Branch-and-Price Algorithm}
\label{sec:bp_impl}
While the Dantzig-Wolfe decomposition provides a tight lower bound at the root node, the RMP solution is typically fractional. To recover an optimal integer solution for the DCMST, column generation is embedded within a Branch-and-Bound framework, yielding a Branch-and-Price algorithm in Appendix \ref{BP}.

\medskip
\noindent 
At each node of the search tree, we solve the LP relaxation of the RMP via column generation until no column with negative reduced cost remains. Let $\lambda_t^*$ denote the optimal multipliers at the current node. We recover the original edge variables by aggregating over all selected trees:
\begin{equation*}
    \bar{x}_e = \sum_{t \in \mathcal{T}:\, e \in t} \lambda_t^* \qquad \forall\, e \in E.
\end{equation*}
If $\bar{x}_e \in \{0,1\}$ for all $e \in E$, the current solution is integral and represents a valid DCMST; if it improves the incumbent upper bound $UB$, we update $x^* \gets \bar{x}$ and $UB \gets Z_{\text{LP}}$. If any edge satisfies $0 < \bar{x}_{e^*} < 1$, the solution is fractional and we must branch.

\medskip
\noindent 
Branching directly on the master variables $\lambda_t$ is impractical in network design, so we branch instead on the original edge variables. We select a fractional edge $e^* = (u,v)$ and create two child nodes. Let $E_0$ and $E_1$ denote the sets of edges fixed to $0$ and $1$ at any given node.
\begin{itemize}
    \item \textbf{Left child} ($x_{e^*} = 0$, $e^* \in E_0$): All columns $t$ containing $e^*$ are removed from the RMP. In the subproblem, $e^*$ is excluded by setting $\bar{w}_{e^*} \gets +\infty$, so Kruskal's algorithm never selects it.
    \item \textbf{Right child} ($x_{e^*} = 1$,\ $e^* \in E_1$): All columns $t$ not containing $e^*$ are removed from the RMP. In the subproblem, $e^*$ is forced into every generated tree by initialising Kruskal's algorithm with $E_1$ pre-selected (merging their components), or equivalently by setting $\bar{w}_{e^*} \gets -\infty$.
\end{itemize}
Remaining valid columns are inherited by each child node as a warm start. These modifications guarantee that column generation only produces trees consistent with the current branching decisions. The algorithm recurses, pruning any node whose LP bound meets or exceeds $UB$, until the optimal integer DCMST is found.

\subsection{Heuristics and Feasibility Recovery}

Because every node must have warm-start column generation with at least one feasible spanning tree, a three-level heuristic cascade is applied both at the root node and at each child node when inherited columns are filtered away.

\begin{enumerate}
  \item \textbf{LP-Guided Greedy:}  Edges are sorted by a composite score that biases towards edges with high LP-solution value $\bar{x}_e$ and low weight: $\text{score}_e = w_e - M \cdot \bar{x}_e$.  When no LP solution is available, edges are sorted by weight alone, recovering a standard Kruskal greedy heuristic.  A Union-Find structure enforces acyclicity and degree constraints are checked before each edge is added.

  \item \textbf{Nearest-Neighbour Heuristic (NNH):}  Edges are sorted in ascending order of the combined remaining-degree sum $(d_u - \deg(u)) + (d_v - \deg(v))$, with weight as a tie-breaker, so that the most constrained vertices are served first.  Edges that would create a cycle or violate a degree bound are skipped.  Any edges in $E_1$ are inserted unconditionally at the start; $E_0$ edges are excluded from consideration.

  \item \textbf{Feasible-Degree-Tree Heuristic:}  At each step the edge $(u,v) \notin E_0$ that minimises the pair $(\text{rem}[u] + \text{rem}[v],\; w_{uv})$ lexicographically—where $\text{rem}[v]$ is the remaining degree budget of $v$—is selected, provided it does not create a cycle.  This strategy prioritises closing off the most degree-constrained vertices and tends to be more robust than NNH when many tight branching constraints are active.
\end{enumerate}

\medskip
\noindent \textbf{MIP feasibility fallback:}  If all three greedy heuristics fail to construct a valid spanning tree, a small Gurobi feasibility model is invoked with a time limit. If this model also fails, the node is proved infeasible and pruned.

\medskip
\noindent After the heuristic cascade, a primal solution update is attempted at every fractional node: the LP-guided greedy is re-run with the current $\bar{x}$ values to round the fractional solution.  Whenever this produces a spanning tree whose cost improves the incumbent, the upper bound is updated and the priority queue is pruned.



\section{Computational Experiments}

\subsection{Experimental Setup}

All benchmark instances are complete undirected graphs on $n$ vertices, where $n$ ranges over $\{20, \dots, 2000\}$.
Edge weights are drawn independently and uniformly at random from the integer interval $[1,\, 2n]$, so the expected weight scale grows with the problem size.
The degree bound is set uniformly across all vertices: $d_v = d$ for all $v \in V$, with $d \in \{2, 3, 4\}$.
For each $(n,d)$ combination, ten independent instances are generated using distinct random seeds $s \in \{0, 1, \ldots, 9\}$, yielding a benchmark of 480 instances in total.


\medskip
\noindent All experiments are conducted on an Intel(R) Core(TM) i5-1021OU CPU @ 1.60 GHZ $\times$ 8, using Gurobi 12.0.1 as the LP and MIP solver throughout.
Each solver is given a wall-clock time limit of 300 seconds and the process is monitored against an 8\,GB memory ceiling; a solve is terminated and labelled accordingly if either resource limit is breached.
An early-stopping rule is applied per method and per degree value: once every seed for a given instance size $n$ exhausts the time or memory budget, all larger sizes are skipped for that method and degree combination. The following four approach are compared.

\begin{enumerate}

  \item \textit{Multi-Commodity Flow:} This method yields $O(n^2)$ additional flow variables and the tightest LP relaxation among the three compact formulations, at the cost of a substantially larger model.
  
  \item \textit{Single-commodity flow:} Classical compact MIP formulation is given to the Gurobi.

  \item \textit{Branch-and-Price:} The implementation is described in detail in Section~\ref{sec:bp_impl}. The Branch-and-Price algorithm is executed in parallel in terms of nodes in branching tree.

    \item \textit{Branch-and-Cut:}  The non-compact formulation in section \ref{sec:non-compact} is passed to Gurobi as an integer program.  Subtour elimination constraints are not added up front; instead, they are separated dynamically as lazy cuts inside a Gurobi callback, which is the standard state-of-practice for this family of formulations.
\end{enumerate}

\subsection{Results}

Table~\ref{tab:addlabel} reports, for each combination of degree bound $d$ and instance size $n$, the average solve time over the ten seeds for which a result was obtained, the number of instances solved to proven optimality within the time limit (O), the number that hit the time limit (T), and the number that exceeded the memory ceiling (M).

\begin{table}[h]
\scriptsize
\setlength{\tabcolsep}{5pt}
  \centering
  \caption{Degree Constraint MST Run Results}
    \begin{tabular}{cccccccccccccccccc}
    \toprule
    \multicolumn{1}{c}{} &       & \multicolumn{4}{c}{\textbf{Multi Comodity Flow}} & \multicolumn{4}{c}{\textbf{Single Comodity Flow}} & \multicolumn{4}{c}{\textbf{Branch \& Price}} & \multicolumn{4}{c}{\textbf{Branch \& Cut}} \\
    \cmidrule(l{3pt}r{3pt}){3-6} \cmidrule(l{3pt}r{3pt}){7-10} \cmidrule(l{3pt}r{3pt}){11-14} \cmidrule(l{3pt}r{3pt}){15-18}
    \multicolumn{1}{c}{\textbf{d}} & \textbf{n} & \textbf{ Time(s)} & \textbf{O} & \textbf{T} & \textbf{M} & \textbf{ Time(s)} & \textbf{O} & \textbf{T} & \textbf{M} & \textbf{ Time(s)} & \textbf{O} & \textbf{T} & \textbf{M} & \textbf{ Time(s)} & \textbf{O} & \textbf{T} & \multicolumn{1}{c}{\textbf{M}} \\
    \midrule
    \multirow{15}[2]{*}{\textbf{2}} & \textbf{20} & 0.244 & 10    & -     & -     & 0.033 & 10    & -     & -     & 0.877 & 10    & -     & -     & 0.015 & 10    & -     & - \\
          & \textbf{50} & 51.212 & 10    & -     & -     & 0.621 & 10    & -     & -     & 36.365 & 10    & -     & -     & 0.041 & 10    & -     & - \\
          & \textbf{80} & 300   & -     & 10    & -     & 3.081 & 10    & -     & -     & 300   & -     & 10    & -     & 0.062 & 10    & -     & - \\
          & \textbf{100} & -     & -     & -     & -     & 6.996 & 10    & -     & -     & -     & -     & -     & -     & 0.12  & 10    & -     & - \\
          & \textbf{150} & -     & -     & -     & -     & 27.356 & 10    & -     & -     & -     & -     & -     & -     & 0.5   & 10    & -     & - \\
          & \textbf{200} & -     & -     & -     & -     & 75.937 & 10    & -     & -     & -     & -     & -     & -     & 0.535 & 10    & -     & - \\
          & \textbf{300} & -     & -     & -     & -     & 256.323 & 7     & 3     & -     & -     & -     & -     & -     & 1.211 & 10    & -     & - \\
          & \textbf{400} & -     & -     & -     & -     & 90.365 & 10    & -     & -     & -     & -     & -     & -     & 3.119 & 10    & -     & - \\
          & \textbf{500} & -     & -     & -     & -     & 167.789 & 7     & 3     & -     & -     & -     & -     & -     & 5.88  & 10    & -     & - \\
          & \textbf{600} & -     & -     & -     & -     & 210.246 & 5     & 5     & -     & -     & -     & -     & -     & 8.814 & 10    & -     & - \\
          & \textbf{700} & -     & -     & -     & -     & 273.864 & 3     & 7     & -     & -     & -     & -     & -     & 12.38 & 10    & -     & - \\
          & \textbf{800} & -     & -     & -     & -     & 300   & -     & 10    & -     & -     & -     & -     & -     & 17.482 & 10    & -     & - \\
          & \textbf{1000} & -     & -     & -     & -     & -     & -     & -     & -     & -     & -     & -     & -     & 28.954 & 10    & -     & - \\
          & \textbf{1200} & -     & -     & -     & -     & -     & -     & -     & -     & -     & -     & -     & -     & 53.513 & 10    & -     & - \\
          & \textbf{1500} & -     & -     & -     & -     & -     & -     & -     & -     & -     & -     & -     & -     & 68.035 & 10    & -     & - \\
          & \textbf{2000} & -     & -     & -     & -     & -     & -     & -     & -     & -     & -     & -     & -     & 147.427 & 8     & -     & 2 \\
    \midrule
    \multirow{8}[2]{*}{\textbf{3}} & \textbf{20} & 2.204 & 10    & -     & -     & 0.035 & 10    & -     & -     & 4.656 & 10    & -     & -     & 0.018 & 10    & -     & - \\
          & \textbf{50} & 157.079 & 10    & -     & -     & 0.52  & 10    & -     & -     & 5.441 & 10    & -     & -     & 0.079 & 10    & -     & - \\
          & \textbf{80} & 300   & -     & 10    & -     & 1.799 & 10    & -     & -     & 7.245 & 10    & -     & -     & 0.248 & 10    & -     & - \\
          & \textbf{100} & -     & -     & -     & -     & 14.316 & 10    & -     & -     & 17.493 & 10    & -     & -     & 1.045 & 10    & -     & - \\
          & \textbf{150} & -     & -     & -     & -     & 31.034 & 10    & -     & -     & 80.606 & 10    & -     & -     & 7.092 & 10    & -     & - \\
          & \textbf{200} & -     & -     & -     & -     & 119.469 & 8     & 2     & -     & 225.141 & 4     & 6     & -     & 34.05 & 10    & -     & - \\
          & \textbf{300} & -     & -     & -     & -     & 234.05 & 4     & 6     & -     & 300   & -     & 10    & -     & 206.093 & 5     & 5     & - \\
          & \textbf{400} & -     & -     & -     & -     & 300   & -     & 10    & -     & -     & -     & -     & -     & 273.912 & 2     & 8     & - \\
          & \textbf{500} & -     & -     & -     & -     & -     & -     & -     & -     & -     & -     & -     & -     & 300   & -     & 10    & - \\
    \midrule
    \multirow{10}[2]{*}{\textbf{4}} & \textbf{20} & 1.273 & 10    & -     & -     & 0.334 & 10    & -     & -     & 4.687 & 10    & -     & -     & 0.048 & 10    & -     & - \\
          & \textbf{50} & 141.033 & 9     & 1     & -     & 6.791 & 10    & -     & -     & 4.464 & 10    & -     & -     & 0.96  & 10    & -     & - \\
          & \textbf{80} & 300   & -     & 10    & -     & 15.569 & 10    & -     & -     & 5.106 & 10    & -     & -     & 2.676 & 10    & -     & - \\
          & \textbf{100} & -     & -     & -     & -     & 99.77 & 8     & 2     & -     & 5.24  & 10    & -     & -     & 9.113 & 10    & -     & - \\
          & \textbf{150} & -     & -     & -     & -     & 90.878 & 8     & 2     & -     & 8.22  & 10    & -     & -     & 40.166 & 10    & -     & - \\
          & \textbf{200} & -     & -     & -     & -     & 212.092 & 6     & 4     & -     & 14.021 & 10    & -     & -     & 160.268 & 9     & 1     & - \\
          & \textbf{300} & -     & -     & -     & -     & 292.48 & 1     & 9     & -     & 28.564 & 10    & -     & -     & 292.637 & 1     & 9     & - \\
          & \textbf{400} & -     & -     & -     & -     & 300   & -     & 10    & -     & 88.642 & 10    & -     & -     & 300   & -     & 10    & - \\
          & \textbf{500} & -     & -     & -     & -     & -     & -     & -     & -     & 257.011 & 3     & 7     & -     & -     & -     & -     & - \\
          & \textbf{600} & -     & -     & -     & -     & -     & -     & -     & -     & 297.148 & 1     & 9     & -     & -     & -     & -     & - \\
          & \textbf{700} & -     & -     & -     & -     & -     & -     & -     & -     & 300   & -     & 10    & -     & -     & -     & -     & - \\
    \bottomrule
    \end{tabular}%
  \label{tab:addlabel}%
\end{table}%

For $d = 2$, the lazy-cut formulation solves all ten instances at every size up to $n = 1500$ and eight of ten at $n = 2000$ (two instances exceed the memory ceiling at that scale), with runtimes that grow smoothly from under a second to approximately 150 seconds. This is the only method that makes meaningful progress beyond $n = 100$ under the harshest degree restriction. Across all three degree values, MCF hits the time limit at $n = 80$ for every seed, and the early-stopping rule eliminates it from larger instances entirely. The single-commodity flow formulation scales further than MCF—reaching $n = 700$ for $d = 2$ and $n = 300$ for $d = 3$ before consistently timing out—but cannot compete with the lazy-cut approach on any degree setting. For $d = 2$, the B\&P algorithm stalls at $n = 80$, performing no better than MCF. At $d = 3$, it scales to $n = 200$ (four of ten optimal) and $n = 300$ (all timing out), roughly matching Integer Flow. The most striking result occurs at $d = 4$: B\&P solves all ten instances at every size up to $n = 400$ with a mean time of 89 seconds, continues to find feasible solutions at $n = 500$ (three of ten) and $n = 600$ (one of ten), and surpasses the lazy-cut formulation, which exhausts its time budget completely at $n = 400$ for $d = 4$. A cross-cutting observation is that increasing $d$ uniformly makes the problem easier for B\&P and Integer Flow but harder for the lazy-cut formulation.

\subsection{Discussion}

The competitive advantage of the lazy-cut approach here stems from the compactness of the DFJ model combined with Gurobi's presolve and cutting-plane machinery, which are highly effective when the feasible region is small and well-structured by tight degree constraints. The $O(n^2)$ flow variable count inflates the LP relaxation to a size that overwhelms both the LP solver and memory, irrespective of how tight the LP bound itself may be. Although MCF provides the tightest LP relaxation among the compact formulations in theory, the practical overhead of a model with $O(n^3)$ non-zeros renders it uncompetitive beyond small instances. Its runtime exhibits non-monotone behaviour (e.g., $n = 400$ is solved faster than $n = 300$ for $d = 2$), which reflects the variance in branch-and-bound tree size across random seeds rather than a structural anomaly. The results confirm the well-known sensitivity of column-generation methods to the degree restriction. This reversal is consistent with the observation that looser degree constraints enlarge the feasible region, producing denser pricing subproblem solutions and tighter column-generation bounds that prune the branch-and-bound tree more aggressively. For Branch \& Cut, $d = 4$ introduces substantially more feasible spanning trees, which means more subtours must be separated before the LP relaxation tightens, inflating both the number of lazy cuts and the branch-and-bound tree. For B\&P, the opposite effect dominates: more feasible trees mean the pricing subproblem finds high-quality columns quickly, yielding tighter LP bounds at the root and reducing the depth of branching required.

% \medskip
% \noindent 
% For $d = 2$, the lazy-cut formulation solves all ten instances at every size up to $n = 1500$ and eight of ten at $n = 2000$ (two instances exceed the memory ceiling at that scale), with runtimes that grow smoothly from under a second to approximately 150 seconds.  This is the only method that makes meaningful progress beyond $n = 100$ under the harshest degree restriction.  The competitive advantage of the lazy-cut approach here stems from the compactness of the DFJ model combined with Gurobi's presolve and cutting-plane machinery, which are highly effective when the feasible region is small and well-structured by tight degree constraints.

% \medskip
% \noindent 
% Across all three degree values, MCF hits the time limit at $n = 80$ for every seed, and the early-stopping rule eliminates it from larger instances entirely.  The $O(n^2)$ flow variable count inflates the LP relaxation to a size that overwhelms both the LP solver and memory, irrespective of how tight the LP bound itself may be.  Although MCF provides the tightest LP relaxation among the compact formulations in theory, the practical overhead of a model with $O(n^3)$ non-zeros renders it uncompetitive beyond small instances.

% \medskip
% \noindent
% The single-commodity flow formulation scales further than MCF—reaching $n = 700$ for $d = 2$ and $n = 300$ for $d = 3$ before consistently timing out—but cannot compete with the lazy-cut approach on any degree setting.  Its runtime exhibits non-monotone behaviour (e.g., $n = 400$ is solved faster than $n = 300$ for $d = 2$), which reflects the variance in branch-and-bound tree size across random seeds rather than a structural anomaly.

% \medskip
% \noindent 
% The results confirm the well-known sensitivity of column-generation methods to the degree restriction.  For $d = 2$, the B\&P algorithm stalls at $n = 80$, performing no better than MCF.  At $d = 3$, it scales to $n = 200$ (four of ten optimal) and $n = 300$ (all timing out), roughly matching Integer Flow.  The most striking result occurs at $d = 4$: B\&P solves all ten instances at every size up to $n = 400$ with a mean time of 89 seconds, continues to find feasible solutions at $n = 500$ (three of ten) and $n = 600$ (one of ten), and surpasses the lazy-cut formulation, which exhausts its time budget completely at $n = 400$ for $d = 4$.  This reversal is consistent with the observation that looser degree constraints enlarge the feasible region, producing denser pricing subproblem solutions and tighter column-generation bounds that prune the branch-and-bound tree more aggressively.

% \medskip
% \noindent 
% A cross-cutting observation is that increasing $d$ uniformly makes the problem easier for B\&P and Integer Flow but harder for the lazy-cut formulation.  For Branch \& Cut, $d = 4$ introduces substantially more feasible spanning trees, which means more subtours must be separated before the LP relaxation tightens, inflating both the number of lazy cuts and the branch-and-bound tree.  For B\&P, the opposite effect dominates: more feasible trees mean the pricing subproblem finds high-quality columns quickly, yielding tighter LP bounds at the root and reducing the depth of branching required.  This complementarity suggests that, in practice, a portfolio approach—selecting the solver based on the degree parameter—would yield the best overall scalability.


\newpage
\appendix
\section{B\&P Algorithm}
\label{BP}
\begin{algorithm}
\caption{Branch-and-Price for Degree-Constrained Minimum Spanning Tree}
\begin{algorithmic}[1]
\footnotesize

\Require Graph $G = (V, E)$, edge weights $w_e$, degree bounds $d_v$
\Ensure Optimal solution $x^*$ and optimal cost $UB$

\Statex
\Statex \textbf{Initialization}
\State $\mathcal{L} \gets \{\text{Root}\}$, \quad $UB \gets \infty$, \quad $x^* \gets \emptyset$
\State For Root: $E_0 \gets \emptyset$, $E_1 \gets \emptyset$

\Statex
\Statex \textbf{Main Loop}
\While{$\mathcal{L} \neq \emptyset$}
    \State Pop node from $\mathcal{L}$ with its sets $E_0$, $E_1$
    \State Remove from $\mathcal{T}$ any column $t$ where $t \cap E_0 \neq \emptyset$ or $E_1 \not\subseteq t$

\Statex
    \Statex \quad \textit{--- Column Generation ---}
    \Repeat
        \State Solve LP relaxation of RMP over current columns $\mathcal{T}$; retrieve dual variables
        \State For pricing: set $\bar{w}_e \gets +\infty$ for $e \in E_0$, and $\bar{w}_e \gets -\infty$ for $e \in E_1$
        \State Solve MST pricing subproblem with modified weights $\bar{w}_e$
        \If{a tree with negative reduced cost is found}
            \State Add it as a new column to $\mathcal{T}$
        \EndIf
    \Until{no column with negative reduced cost exists}
    \State Let $Z_{\text{LP}}$ be the LP objective value and $\lambda^*$ the optimal multipliers

\Statex
    \Statex \quad \textit{--- Branch and Bound ---}
    \If{$Z_{\text{LP}} \geq UB$}
        \State \textbf{Prune} by bound; \textbf{continue}
    \EndIf
    \State Compute $\bar{x}_e \gets \sum_{t \in \mathcal{T}:\, e \in t} \lambda_t^*$ \quad $\forall e \in E$
    \If{$\bar{x}_e \in \{0,1\}$ for all $e \in E$}
        \If{$Z_{\text{LP}} < UB$}
            \State $UB \gets Z_{\text{LP}}$, \quad $x^* \gets \bar{x}$
        \EndIf
        \State \textbf{Prune} (integral); \textbf{continue}
    \EndIf

    \State Select fractional edge $e^* \in E$ with $0 < \bar{x}_{e^*} < 1$
    \State \textbf{Left child} (exclude $e^*$): \enspace $E_0' \gets E_0 \cup \{e^*\}$, \enspace $E_1' \gets E_1$; \enspace add to $\mathcal{L}$
    \State \textbf{Right child} (include $e^*$): \enspace $E_0' \gets E_0$, \enspace $E_1' \gets E_1 \cup \{e^*\}$; \enspace add to $\mathcal{L}$

\EndWhile 
\Statex \textbf{Termination}
\State \Return $x^*$, $UB$

\end{algorithmic}
\end{algorithm}

\newpage
\begin{thebibliography}{99}

\bibitem[Behle et al.(2007)]{Behle2007}
Behle, M., J\"unger, M., \& Liers, F. (2007). 
\newblock A primal branch-and-cut algorithm for the degree-constrained minimum spanning tree problem. 
\newblock \emph{Experimental Algorithms}, 379--392.

\bibitem[Bicalho et al.(2016)]{Bicalho2016}
Bicalho, L. H., da Cunha, A. S., \& Lucena, A. (2016). 
\newblock Branch-and-cut-and-price algorithms for the degree constrained minimum spanning tree problem. 
\newblock \emph{Computational Optimization and Applications}, 63(1), 227--253.

\bibitem[Caccetta \& Hill(2001)]{Caccetta2001}
Caccetta, L., \& Hill, S. P. (2001). 
\newblock A branch and cut method for the degree-constrained minimum spanning tree problem. 
\newblock \emph{Networks}, 37(2), 74--83.

\bibitem[Fisher(1981)]{Fisher1981}
Fisher, M. L. (1981). 
\newblock The Lagrangian relaxation method for solving integer programming problems. 
\newblock \emph{Management Science}, 27(1), 1--18.

\bibitem[Gamarnik et al.(2005)]{Gamarnik2005}
Gamarnik, D., Lewenstein, M., \& Sviridenko, M. (2005). 
\newblock An improved upper bound for the TSP in cubic 3-edge-connected graphs. 
\newblock \emph{Operations Research Letters}, 33(5), 467--474.

\bibitem[Gouveia et al.(2011)]{Gouveia2011}
Gouveia, L., Leitner, M., \& Ljubi\'c, I. (2011). 
\newblock On the topological design of networks with diameter constraints. 
\newblock \emph{INFORMS Journal on Computing}, 23(1), 17--36.

\bibitem[Knowles \& Corne(1999)]{Knowles1999}
Knowles, J., \& Corne, D. (1999). 
\newblock A new evolutionary approach to the degree-constrained minimum spanning tree problem. 
\newblock \emph{IEEE Transactions on Evolutionary Computation}, 4(2), 125--134.

\bibitem[Volgenant(1989)]{Volgenant1989}
Volgenant, A. (1989). 
\newblock A Lagrangean approach to the degree-constrained minimum spanning tree problem. 
\newblock \emph{European Journal of Operational Research}, 39(3), 325--331.

\end{thebibliography}
\end{document}

